% SPDX-License-Identifier: CC-BY-4.0
% Copyright (c) 2026 Martin Schröder <info@swedishembedded.com>

\section{The System}\label{sec:system}

Splinter is a learning agent built on two dependencies: an agent runtime that runs model
calls and tools, and a training and inference engine (called \emph{brain} here) that
performs all model computation. Splinter owns the learning campaign: what to teach, which
evidence qualifies as training material, how a candidate is measured, and whether it is
released. A run is one recorded pipeline of stages (\cref{fig:pipeline}); every stage stores
its output under a content address so that any stage can be rerun or inspected alone.

\begin{figure}[t]
\centering
\resizebox{\textwidth}{!}{%
\begin{tikzpicture}[
  node distance=5mm and 6mm,
  box/.style={draw, rounded corners=2pt, minimum height=9mm, minimum width=19mm, align=center, font=\scriptsize, fill=blue!6},
  gate/.style={draw, rounded corners=2pt, minimum height=9mm, minimum width=19mm, align=center, font=\scriptsize, fill=orange!14},
  arr/.style={-{Stealth[length=2mm]}, thick}]
\node[box] (plan) {plan};
\node[box, right=of plan] (src) {source\\capture};
\node[box, right=of src] (res) {reserve\\exam families};
\node[box, right=of res] (tasks) {task\\generation};
\node[box, right=of tasks] (solve) {solve,\\verify};
\node[box, below=of plan] (teach) {teach\\(open book)};
\node[box, right=of teach] (author) {author\\(verbatim)};
\node[box, right=of author] (data) {dataset\\build};
\node[box, right=of data] (train) {LoRA\\fine-tune};
\node[gate, right=of train] (exam) {exam on\\held-out\\families};
\node[gate, below=of exam] (rel) {release\\gate};
\draw[arr] (plan) -- (src); \draw[arr] (src) -- (res); \draw[arr] (res) -- (tasks);
\draw[arr] (tasks) -- (solve);
\draw[arr] (solve.south) -- ++(0,-3mm) -| ($(plan.south)+(0,-3mm)$) -- (teach.north);
\draw[arr] (teach) -- (author); \draw[arr] (author) -- (data); \draw[arr] (data) -- (train);
\draw[arr] (train) -- (exam); \draw[arr] (exam) -- (rel);
\end{tikzpicture}}
\caption{The stages of one run, in order (the two orange stages measure). The reserve stage
precedes every generating stage, so no exam text is ever read by anything that makes
training data. Rehearsal and variant stages are omitted for clarity.}
\label{fig:pipeline}
\end{figure}

\paragraph{Models and hardware.}
Policy, teacher, task generator, planner and critic are all Qwen3-8B in non-thinking mode.
The judge is Qwen3-14B, a separate model loaded on its own. Two Tesla P40 cards (24~GB each)
are used, one job per card; the engine runs on them through Vulkan, and the policy base is
held in bf16.
Every job is launched through a wrapper that pins the card and stops the run if it
touches a CPU backend or software adapter, so a measurement can never silently come from
the wrong device.

\subsection{Stages}

\paragraph{Plan.}
A planner model reads a survey of the materials and chooses what to teach from a menu
of task kinds (recall, the author's own advice to a correspondent, explanation, conversation,
restoration of corrupted passages). Each choice is held to the survey by code: a kind
the sources cannot support is refused.

\paragraph{Source capture and reservation.}
The directory is captured as a source of text \emph{parts} (the Jefferson materials hold
1,629 parts; the Adams materials 307), split into sections that tasks cite by byte range.
Before any task is generated, the \emph{reserve} stage picks the families of the final test
and of a dev suite by a stable hash of the family name and removes the text of every edition
of each from every source later stages read (\cref{sec:split}).

\paragraph{Task generation.}
A generator model proposes tasks from sections; code admits or refuses each one. A task
must be grounded in its section, stand on its own without the source in front of the
student, name its subject, quote its source at a run of eight words rather than
paraphrase it, and not contradict another task. In the Jefferson run described below,
the generator was given 6,477~s and covered 19 of 1,629 parts; at that pace covering every
part would have taken 555,322~s, so only a sliver of the corpus is ever turned into tasks.
For the persona runs the generator is told the person wrote the sources.

\paragraph{Solve and verify.}
The policy attempts each task closed-book; tasks it always solves are dropped. A teacher
(here the same Qwen3-8B, shown the source passage) answers the rest open-book. Verifiers
are ranked by strength: executable checks, a \emph{stated} check that accepts an answer
stating the reference, an agreement check, and a calibrated judge. The strongest verdict
decides, and a judge alone never makes a training record. The student's training record
keeps the instruction and the verified answer but never the passage.

\paragraph{Author channel and variants.}
For a person, a second kind of record asks for a passage of the person's own text (at most 85
words) and keeps it only if a fit judge says the passage is a natural reply to the request.
Separately, the generator rewrites each kept task in other words; these \emph{variants} are
never trained on.

\paragraph{Dataset views.}
Records are projections of stored experience into training format. Three views matter here.
The \emph{dialogue} view holds verified teacher answers under the persona system prompt
(``You are Thomas Jefferson. Answer as Thomas Jefferson would: in the first person, in
Thomas Jefferson's own voice, from Thomas Jefferson's own experience and opinions. Answer the
user's request directly and concisely.''). The \emph{voice} view is built by code with no model
in the loop: the whole corpus of the author's text in chunks of 250 to 650 words, whole
sentences, closed at paragraph ends, one print of each text, no family over a tenth of the
tokens, half of them asked under the identity line instead of the persona prompt, each asked
for by a request built by code from its recipient, year and opening. In a later arm the
request is a 20--60 word description of the chunk's content written by the generator and
refused if more than 35\% of its words fall in runs of four or more words copied from the
passage. The
\emph{rehearsal} view holds the base model's own greedy answers to general tasks (sums and
format requests built by code with references, and general requests the base writes), kept
unless code decides them wrong. An optional abstention share trains the author to say the
writings do not settle a question; it is kept under half of the records, because a model
taught to stop that often stops answering.

\paragraph{Training.}
A candidate is a rank-8 LoRA adapter on all seven projections of every layer
(q, k, v, o, gate, up, down), trained by supervised fine-tuning in no-think mode
with loss only on the answer tokens and, in every run reported here, eight records per
step.
The current recipe is one learning rate for both entry commands (\code{train} and
\code{learn}, 2e-4), an alpha of twice the rank, no weight decay, AdamW with gradient
clipping at 1.0, a warmup of a twentieth of the steps, a cosine decay to a floor of a tenth
of the peak over the planned steps, and a step that is the mean over its supervised tokens,
with every record drawn once per epoch. A share of whole training families (a tenth of the
records, over at least three families where the data have them) is held back as a
\emph{monitoring} set scored at evaluations a sixteenth of the budget apart unless named
(every 15 steps for e4 and e1); the run keeps the adapter
of the best monitoring loss, and after four evaluations without improvement it lowers the
rate linearly to the floor over a tenth of the planned steps, still evaluating, and stops.
Of the candidates in \cref{sec:results}, e4 and e1 were trained with this recipe and all
others with the earlier one (\cref{sec:training-defects}).
Runs can keep the adapter of every evaluation so that a checkpoint may later be chosen on a
dev suite by generation quality rather than loss.

\paragraph{Exam and release gate.}
The exam (\cref{sec:eval}) measures whether the candidate generalises to held-out families.
The release gate has four checks: improvement over the champion on held-out and paraphrased
tasks (a sign test at $\alpha=0.05$), retention on earlier releases' tasks (bound 0.05),
a frozen \emph{anchor} suite of general tasks (110 items: 70 recall, 28 arithmetic, 12
format-following) with a drop bound of 0.02, and a serve check that plain \code{brain serve}
with the adapter gives the same answers on eight sampled tasks (at least 75\% agreement).
The improvement decision is the family-level sign test; task-level counts are reported
beside it. A release is an immutable adapter with a manifest of every number the gate
measured.

\subsection{Residency and memory on 24~GB cards}\label{sec:residency}

A bf16 Qwen3-8B occupies about 16.4~GB, and training at rank 8 with 1,024-token rows reached
a sampled peak of 21,407~MiB on one card. The 14B judge is a separate base that is not resident beside it. The
design consequence is that a device holds one base at a time and the work is staged so
that bases swap as seldom as possible: the policy, the candidate and the champion are one
base with different adapters attached before each generation (attached, never folded in,
because folding an adapter into a base forces a re-read of the checkpoint on every
switch); the gate and the exam answer \emph{every} task of a suite with every arm before the
judge loads; every resident base is released before a fine-tune or a \code{brain serve}
process loads its own copy; the serve check judges its answers only after the server has
exited; and the voice score loads one arm at a time and gives the device back between arms.
These rules came from failures: an out-of-memory fault on a card is not recoverable in
process (the engine reports a hard device fault and leaks the device rather than destroy
inconsistent driver state, so the run is lost), and in the earliest gate runs a judge that loaded while the server still held the
card ran out of memory, left every served answer ungraded, and the check counted ungraded
as ``different''. A served task with no answer is now ``not measured'', never a disagreement.
